% Два направления методов штрафов: внешний штраф и барьерная функция.
\begin{figure}[htbp]
\centering
\begin{tikzpicture}[scale=1.0, >={Stealth[length=2mm]}]

% ---------- Левая панель: внешний штраф ----------
\begin{scope}
  \draw[->] (-0.3,0) -- (4.6,0) node[right, font=\footnotesize] {$x$};
  \draw[->] (0,-0.3) -- (0,3.4) node[above, font=\footnotesize] {$\Phi_r$};

  % границы допустимой области
  \draw[gray, dashed] (1.0,-0.15) -- (1.0,3.2);
  \draw[gray, dashed] (3.2,-0.15) -- (3.2,3.2);
  \node[font=\scriptsize, anchor=north] at (1.0,-0.2) {$a$};
  \node[font=\scriptsize, anchor=north] at (3.2,-0.2) {$b$};

  % внутри области: сама целевая функция
  \draw[blue!55!black, very thick, domain=1.0:3.2, samples=60, smooth]
    plot (\x, {0.55*(\x-2.1)*(\x-2.1) + 0.5});
  % снаружи: резкий рост штрафа
  \draw[orange!75!black, very thick, domain=0.25:1.0, samples=60, smooth]
    plot (\x, {0.55*(\x-2.1)*(\x-2.1) + 0.5 + 9.0*(1.0-\x)*(1.0-\x)});
  \draw[orange!75!black, very thick, domain=3.2:3.75, samples=60, smooth]
    plot (\x, {0.55*(\x-2.1)*(\x-2.1) + 0.5 + 9.0*(\x-3.2)*(\x-3.2)});

  \node[blue!55!black, font=\scriptsize, anchor=south] at (2.1,1.0)
    {$\equiv \varphi_0$};
  \node[font=\footnotesize, anchor=north] at (2.1,-0.75) {внешний штраф};
\end{scope}

% ---------- Правая панель: барьер ----------
\begin{scope}[xshift=6.3cm]
  \draw[->] (-0.3,0) -- (4.6,0) node[right, font=\footnotesize] {$x$};
  \draw[->] (0,-0.3) -- (0,3.4) node[above, font=\footnotesize] {$\Phi_r$};

  \draw[gray, dashed] (1.0,-0.15) -- (1.0,3.2);
  \draw[gray, dashed] (3.2,-0.15) -- (3.2,3.2);
  \node[font=\scriptsize, anchor=north] at (1.0,-0.2) {$a$};
  \node[font=\scriptsize, anchor=north] at (3.2,-0.2) {$b$};

  % барьер: уходит в +бесконечность у границ изнутри
  \draw[orange!75!black, very thick, domain=1.08:3.12, samples=160, smooth]
    plot (\x, {0.45*(\x-2.1)*(\x-2.1) + 0.45
               + 0.035/(\x-1.0) + 0.035/(3.2-\x)});

  \node[orange!60!black, font=\scriptsize, anchor=west] at (3.25,2.6)
    {$\to +\infty$};
  \node[font=\footnotesize, anchor=north] at (2.1,-0.75)
    {барьер (внутренний штраф)};
\end{scope}

\end{tikzpicture}
\caption{Два направления методов штрафов. Слева внешний штраф: внутри области
$\Phi_r$ совпадает с $\varphi_0$, снаружи резко растёт, и приближения подходят
к $S$ снаружи. Справа барьерная функция: она обращается в $+\infty$ при подходе
к границе изнутри, поэтому точка не может покинуть $S$, но $\varphi_0$
искажается и внутри области — искажение снимается лишь при $r \to 0$.}
\label{fig:penalty-barrier}
\end{figure}
